% Generated draft: requires mathematical and source review.
\documentclass[11pt]{article}
\usepackage[margin=25mm]{geometry}
\usepackage{fontspec}
\setmainfont{Latin Modern Roman}
\usepackage{amsmath,amssymb,mathtools,unicode-math}
\setmathfont{Latin Modern Math}
\usepackage{enumitem,needspace,xurl,hyperref}
\setlength{\parindent}{0pt}
\setlength{\parskip}{5pt}
\setlength{\emergencystretch}{3em}
\newcommand{\N}{\mathbb{N}}
\newcommand{\Z}{\mathbb{Z}}
\newcommand{\Q}{\mathbb{Q}}
\newcommand{\R}{\mathbb{R}}
\newcommand{\C}{\mathbb{C}}
\newcommand{\F}{\mathbb{F}}
\newcommand{\A}{\mathbb{A}}
\newcommand{\PP}{\mathbb{P}}
\newcommand{\E}{\mathbb{E}}
\newcommand{\eps}{\varepsilon}
\newcommand{\dd}{\,\mathrm{d}}
\begin{document}
\section*{UnsolvedMath ID: 30002127}
\textbf{Source problem code:} catalogue entry
\subsection{Definitions and mathematical statement}
A complete Riemannian plane is understood here to be a smooth surface \(P\), without boundary and diffeomorphic to \(\R^2\), equipped with a complete Riemannian metric \(g\). Under the convention that a geodesic means the unparametrized image of a complete nonconstant geodesic, put
\[
\mathcal G(P,g)=\{\gamma(\R):\gamma\colon\R\to P\text{ is a unit-speed geodesic}\}.
\]
The parallel axiom is the assertion
\[
\forall L\in\mathcal G(P,g)\;\forall p\in P\setminus L,\qquad
\#\{M\in\mathcal G(P,g):p\in M\text{ and }M\cap L=\varnothing\}=1.
\]
Thus uniqueness refers to the geodesic image, independently of its parametrization or orientation.

\emph{Question.} If a complete Riemannian plane \((P,g)\) satisfies this parallel axiom, must there exist a Riemannian isometry
\[
F\colon(P,g)\longrightarrow(\R^2,\dd x^2+\dd y^2)?
\]
No curvature condition or absence-of-conjugate-points hypothesis is stated in the supplied problem.

\emph{Interpretive qualification.} The supplied record does not explicitly define ``Riemannian plane'' or specify whether ``geodesic'' means every complete geodesic or only a globally minimizing geodesic line. The formulation above uses all complete geodesic images; it does not impose global minimization. These conventions require confirmation against the source. The registry labels the question open; this is a supplied status, not an independently verified mathematical conclusion.
\subsection{Short English statement}
Must every complete Riemannian surface diffeomorphic to \(\R^2\) be Euclidean if, through each point outside each complete geodesic, there passes exactly one complete geodesic disjoint from it? Here geodesics are counted by their images; the source's intended convention concerning globally minimizing lines requires review.
\subsection{Sources}
\begin{itemize}
\item Supplied local UnsolvedMath registry record, permanent ID \(30002127\), title: \emph{Parallel-Axiom Rigidity of Complete Riemannian Planes}; collection: \emph{Oberwolfach Reports Open Problems}; exact source code: \texttt{catalogue entry}. The supplied proposed year is 2012; no individual proposer is identified.
\item \emph{Geometric Group Theory, Hyperbolic Dynamics and Symplectic Geometry} (2013), \emph{Oberwolfach Reports}, pp.\ 2139--2203, DOI 10.4171/owr/2012/35: \url{https://doi.org/10.4171/owr/2012/35}. This is also listed among the supplied literature sources as ``Oberwolfach Report 35.'' The background describes these reports as documentation of mathematical workshops and their open-problem discussions.
\item Supplied literature source: \emph{Rigidity of complete manifolds without conjugate points}, \url{https://d-nb.info/1123477280/34}.
\item Supplied literature source: \emph{Bangert--Emmerich, area growth and rigidity}, \url{https://archive.ymsc.tsinghua.edu.cn/pacm_download/21/980-JDG-94-3-A1-Bangert.pdf}.
\item Supplied literature assessment, checked on \texttt{2026-08-21}: the record labels the problem open and states that subsequent rigidity literature still identifies this parallel-axiom rigidity question as open. This is supplied dated curation derived from the Oberwolfach material, not an independently verified assessment; time-sensitive status and cited claims require independent verification.
\item The supplied statement status is \texttt{corrected\_verified}. Its verification note reports that the mathematical content was checked against the official Oberwolfach report and that notation, context, and question form were normalized without changing the source-level claim. This draft preserves that reported qualification without claiming an independent source check.
\end{itemize}
\end{document}
